\documentclass[10pt,a4paper]{article}

\usepackage[margin=0.48in]{geometry}
\usepackage[T1]{fontenc}
\usepackage[utf8]{inputenc}
\usepackage[spanish,es-tabla]{babel}
\usepackage{titlesec}
\usepackage{enumitem}
\usepackage[hidelinks]{hyperref}
\usepackage{parskip}

\titleformat{\section}{\normalsize\bfseries\uppercase}{}{0em}{}[\titlerule]
\titlespacing*{\section}{0pt}{6pt}{3pt}
\setlist[itemize]{noitemsep, topsep=0pt, leftmargin=1.1em, itemsep=1pt}
\setlength{\parskip}{1pt}

\begin{document}
\fontsize{9.5pt}{11.8pt}\selectfont

\begin{center}
    {\LARGE\textbf{Emilio Gordillo Esparragoza}}\\[3pt]
    Ingeniero de Datos y Desarrollador Backend\\[4pt]
    \href{mailto:emilio.gordillo.dev@gmail.com}{emilio.gordillo.dev@gmail.com}
    $\cdot$ M\'exico
    $\cdot$ \href{https://github.com/Emilio-Gordillo-Esparragoza}{github.com/Emilio-Gordillo-Esparragoza}
    $\cdot$ \href{https://www.linkedin.com/in/emilio-gordillo-esparragoza-948267397/}{LinkedIn}
\end{center}

\section*{Perfil Profesional}
Aspirante a Ingeniero de Datos y Desarrollador Backend con experiencia en Python, SQL/PostgreSQL y arquitecturas AWS. Construye sistemas productivos de rutas con TypeScript, Supabase y PostgreSQL mientras cursa Ingenier\'ia de Software (UAG). Dise\~na pipelines de datos, automatiza entornos DevOps e implementa patrones serverless (Lambda, SQS, EventBridge, Terraform). Busca contribuir en pipelines ETL/ELT, APIs y servicios backend.

\section*{Habilidades T\'ecnicas}
\textbf{Lenguajes:} Python, SQL, TypeScript/JavaScript, PLpgSQL, Java, Rust\\
\textbf{Datos y An\'alisis:} PostgreSQL, Pandas, NumPy, Apache Spark, scikit-learn, an\'alisis estad\'istico (ANOVA, pruebas t)\\
\textbf{Nube y Backend:} AWS (Lambda, SQS, EventBridge, serverless), APIs REST, Supabase, FastAPI, Flask\\
\textbf{DevOps y Herramientas:} Docker, Terraform, CI/CD (GitHub Actions, Jenkins), Ansible, Nginx, Git

\section*{Experiencia}

\textbf{Desarrollador de Software} \hfill Junio 2026 -- Presente\\
\href{https://www.enrutate.com/}{Enrutate} --- Plataforma de transporte en Culiac\'an y Guasave\\
\begin{itemize}
    \item Desarrolla y mantiene una aplicaci\'on productiva de seguimiento de rutas para Culiac\'an y Guasave con TypeScript/JavaScript, Supabase y PostgreSQL (PLpgSQL).
    \item Implementa actualizaciones en tiempo real y validaciones del servidor para mejorar precisi\'on de rutas y experiencia de usuarios.
\end{itemize}

\section*{Proyectos}

\textbf{\href{https://github.com/Emilio-Gordillo-Esparragoza/thinking-serverless}{thinking-serverless}} --- AWS Serverless y Arquitectura Orientada a Eventos\\
\begin{itemize}
    \item Proyectos desplegables con fan-out EventBridge, flujos SQS/Lambda, colas FIFO/DLQ; infraestructura con Terraform e idempotencia.
\end{itemize}

\textbf{\href{https://github.com/Emilio-Gordillo-Esparragoza/naim-simple}{naim-simple}} --- Pipeline ML para Datos Faltantes (Python)\\
\begin{itemize}
    \item API estilo scikit-learn (fit/predict/save/load) para imputaci\'on NAIM y preprocesamiento reutilizable en datasets tabulares.
\end{itemize}

\textbf{\href{https://github.com/Emilio-Gordillo-Esparragoza/devtool-pack}{devtool-pack}} --- Automatizaci\'on Cross-Platform de Entornos DevOps\\
\begin{itemize}
    \item Paquete Python que instala Terraform, AWS CLI, Docker y kubectl en Windows/macOS/Linux, reduciendo setup de horas a un comando.
\end{itemize}

\textbf{\href{https://github.com/Emilio-Gordillo-Esparragoza/PonyFlow}{PonyFlow}} --- Sistema Multiagente de Generaci\'on de C\'odigo\\
\begin{itemize}
    \item Pipeline multiagente auto-correctivo (LangChain/Ollama) con UI de escritorio React + Tauri (Rust).
\end{itemize}

\section*{Educaci\'on}
\textbf{Universidad Aut\'onoma de Guadalajara (UAG)} \hfill 2026 -- Presente\\
Ingenier\'ia de Software

\section*{Certificaciones}
IBM Data Engineering Professional Certificate $\cdot$
AWS Academy Cloud Foundations $\cdot$
AWS CloudQuest: Cloud Practitioner $\cdot$
AWS Academy Generative AI Foundations $\cdot$
IBM Applied Software Engineering Fundamentals

\section*{Idiomas}
Ingl\'es --- Fluido \quad|\quad Espa\~nol --- Nativo

\end{document}
